\documentclass[acmsmall]{acmart}

%% Rights management
\setcopyright{acmlicensed}
\copyrightyear{2026}
\acmYear{2026}
\acmDOI{XXXXXXX.XXXXXXX}

%% Journal information
\acmJournal{IMWUT}
\acmVolume{10}
\acmNumber{2}
\acmArticle{1}
\acmMonth{6}

%% Bibliography style
\bibliographystyle{ACM-Reference-Format}

\begin{document}

\title{Just a Glimpse: Temporal Wearing as a Design Paradigm for Transient Mixed Reality}

\author{Botao Amber Hu}
\email{botao.hu@cs.ox.ac.uk}
\affiliation{%
  \institution{University of Oxford}
  \city{Oxford}
  \country{United Kingdom}
}

\author{Yilan Tao}
\email{yilan_tao@sfu.ca}
\affiliation{%
  \institution{Simon Fraser University}
  \city{Burnaby}
  \state{BC}
  \country{Canada}
}

\author{Danlin Huang}
\email{danlin@caa.edu.cn}
\affiliation{%
  \institution{China Academy of Art}
  \city{Hangzhou}
  \country{China}
}

\author{Bernhard E. Riecke}
\email{ber1@sfu.ca}
\affiliation{%
  \institution{Simon Fraser University}
  \city{Burnaby}
  \state{BC}
  \country{Canada}
}

\author{Yue Li}
\email{yue.li@xjtlu.edu.cn}
\affiliation{%
  \institution{Xi'an Jiaotong-Liverpool University}
  \city{Suzhou}
  \country{China}
}

\begin{abstract}
Mixed reality devices are overwhelmingly designed for sustained wearing---sessions measured in minutes or hours. Yet many of the most compelling MR encounters in museums, exhibitions, and public spaces last only seconds: a visitor raises a viewer to their eyes, glimpses an augmented world, sets the device down, and returns moments later for another look. We call this pattern \textit{temporal wearing}---an interaction paradigm where the act of donning and doffing a device \textit{is} the interaction, not friction around it. In this paper, we formalize temporal wearing as a distinct category alongside permanent and session wear, proposing a design space defined by body coupling, transition cost, and social legibility. Through a Research through Design process, we developed six non-head-mounted MR prototypes spanning handheld, neck-hung, and stationary form factors, each designed for seconds-long encounters. We deployed three prototypes in a public exhibition (N=XX) and analyzed interaction logs, video observations, and interviews (N=XX). Our findings reveal that lower transition costs encourage repeated re-engagement, that temporal wearing enables socially shared MR experiences, and that visitors develop distinctive temporal wearing behaviors including ``double-takes,'' ``hand-offs,'' and ``orbiting'' patterns. We contribute the temporal wearing concept, a three-part wearing duration taxonomy, empirical evidence from field deployment, and six design principles for temporal wearing devices.
\end{abstract}

\begin{CCSXML}
<ccs2012>
<concept>
<concept_id>10003120.10003121.10003129</concept_id>
<concept_desc>Human-centered computing~Interactive systems and tools</concept_desc>
<concept_significance>500</concept_significance>
</concept>
<concept>
<concept_id>10003120.10003121.10003125.10011752</concept_id>
<concept_desc>Human-centered computing~Haptic devices</concept_desc>
<concept_significance>300</concept_significance>
</concept>
<concept>
<concept_id>10003120.10003121.10003122.10003334</concept_id>
<concept_desc>Human-centered computing~User studies</concept_desc>
<concept_significance>300</concept_significance>
</concept>
</ccs2012>
\end{CCSXML}

\ccsdesc[500]{Human-centered computing~Interactive systems and tools}
\ccsdesc[300]{Human-centered computing~Haptic devices}
\ccsdesc[300]{Human-centered computing~User studies}

\keywords{temporal wearing, mixed reality, wearable computing, transient interaction, exhibition design, social acceptability}

\maketitle

\section{Introduction}

A visitor enters a gallery and notices a small, unassuming device resting on a plinth beside a painting. Curious, she picks it up and raises it to her eyes. Through the viewfinder, the painting comes alive---brushstrokes lift from the canvas, the artist's preparatory sketches shimmer beneath the surface, and a constellation of historical references unfolds in the space around the frame. After a few seconds, she lowers the device and sets it back on the plinth. She moves on to the next room. But something pulls her back. She returns, picks up the viewer again, and this time notices details she missed---a hidden signature, a color shift that reveals a later restoration. She sets it down once more, turns to her companion, and says: ``You have to see this.''

This encounter---brief, voluntary, repeatable, and shareable---represents a form of mixed reality interaction that is both remarkably common and remarkably undertheorized. Museum viewfinders, exhibition stereoscopes, coin-operated binoculars at scenic overlooks, and increasingly, handheld AR experiences at cultural sites all follow the same pattern: a person picks up an optical device, holds it to their body for seconds, experiences an augmented view, and puts it down. The interaction is not a session to be endured but a \textit{glimpse} to be savored.

Yet the dominant approach in mixed reality design assumes sustained use. Head-mounted displays from research prototypes to commercial products---Meta Quest, Apple Vision Pro, Microsoft HoloLens---are engineered for wearing durations measured in minutes to hours. Their design priorities reflect this assumption: ergonomic weight distribution for extended comfort, adjustable straps for secure fit, ventilation to manage heat buildup, and software ecosystems that reward prolonged engagement. Donning and doffing these devices is treated as \textit{friction}---an overhead cost to be minimized so that the ``real'' interaction can begin \cite{knight2006comfort}. The implicit message is clear: the longer you wear it, the better.

This assumption creates a fundamental mismatch when MR encounters are transient. In museums, the average dwell time at an exhibit is 17--30 seconds \cite{vomlehn2001exhibiting}. In exhibitions and public demonstrations, visitors cycle through experiences in minutes. In educational settings, students pass AR devices around for brief, focused encounters. In all of these contexts, the interaction \textit{is} the wearing---not what happens during a sustained session, but the act of bringing a device to the body, glimpsing an augmented reality, and setting it down. Designing these encounters with devices optimized for sustained wear is like designing a doorway with the assumption that people will stand in it for an hour.

The mismatch extends beyond ergonomics into social dynamics. Research on the social acceptability of head-worn devices consistently demonstrates that wearables covering the face provoke discomfort, suspicion, and stigma among bystanders \cite{koelle2020social, koelle2017acceptability}. People wearing HMDs in public spaces appear disconnected, surveilling, or simply odd. These social costs are proportional to wearing duration: an hour in a VR headset at a caf\'{e} reads very differently from a two-second peek through a museum viewfinder. Yet social acceptability research has largely examined continuous wearing scenarios, leaving the temporal dimension unexplored.

We propose \textbf{temporal wearing} as an undertheorized design dimension that reframes this mismatch. Temporal wearing describes encounters where wearing a device for a brief, bounded duration---typically seconds---\textit{constitutes} the interaction rather than enabling it. In temporal wearing, donning and doffing are not overhead costs but productive acts: the don initiates a perceptual shift, the doff concludes it, and the cycle of donning and doffing structures the entire experience. The device is not a tool that happens to be worn; the wearing \textit{is} the tool.

This practice is not new. It has deep historical roots in optical culture. The lorgnette, popular from the 18th century onward, was a pair of lenses on a handle---raised to the eyes for a focused look, then lowered. Opera glasses served the same function in theatrical contexts. The monocle, the hand-held magnifying glass, the stereoscope, the tourist viewfinder: all are temporal wearing devices, designed not for sustained use but for repeated, brief encounters. What is new is not the practice but the \textit{formalization}. Despite centuries of temporal wearing and decades of wearable computing research, no framework has identified wearing duration as a primary design variable or theorized the don/doff cycle as an interaction in its own right. We argue that formalizing this long-standing practice has measurable design consequences: it reveals a design space that existing frameworks overlook and yields empirically grounded principles for a category of wearable interaction that has been practiced but never systematically studied.

Foregrounding wearing duration as a design dimension reveals relationships that existing frameworks miss. Where wearability research asks \textit{where} on the body a device should be placed \cite{gemperle1998wearability, zeagler2017where} and \textit{how} it should be attached, temporal wearing asks \textit{for how long} and \textit{how often}. Where transition research examines how users move between devices or interface states \cite{hubenschmid2025hybrid}, temporal wearing examines how users move between wearing and not-wearing as a \textit{cyclic} process---not a single transition from state A to state B, but a repeating rhythm of engagement and disengagement. And where calm technology envisions information moving between the center and periphery of attention \cite{weiser1996calm}, temporal wearing extends this to the device itself: the device moves between presence and absence, between body-coupled and body-decoupled states.

In this paper, we develop the temporal wearing concept and validate it through design and empirical study. Our contributions are:

\begin{itemize}
    \item[\textbf{C1}] \textbf{Temporal wearing as a design dimension} and a three-part \textbf{wearing duration taxonomy} (permanent, session, and temporal wear) that foregrounds an overlooked but consequential variable in wearable interaction design. Temporal wearing has been practiced for centuries but never formalized; we argue that formalizing it has measurable design consequences.
    \item[\textbf{C2}] A \textbf{design space} for temporal wearing devices defined by three measurable axes---body coupling, transition cost, and social legibility---populated with six prototype instantiations and historical precedents.
    \item[\textbf{C3}] \textbf{Two-phase empirical validation}: a controlled lab study (N=20) comparing temporal wearing prototypes against a session-wear baseline on transition time, re-engagement, social comfort, and cognitive load; and an in-situ deployment at a public exhibition revealing emergent temporal wearing behaviors.
    \item[\textbf{C4}] Six \textbf{design principles} for temporal wearing devices, grounded in empirical findings from both study phases, offering actionable guidance for designing MR encounters measured in seconds rather than sessions.
\end{itemize}

\section{Related Work}

Our work draws on three bodies of literature that, taken together, reveal a gap: mixed reality is increasingly encountered in transient contexts, optical wearing has always been temporal but remains untheorized in wearable computing, and the cyclic nature of temporal interaction has been overlooked by transition frameworks. We review each strand before articulating how temporal wearing unifies them.

\subsection{Transient Mixed Reality in Public Contexts}

Mixed reality, as encountered by the public, is overwhelmingly transient. Museum and gallery deployments, exhibition demonstrations, pop-up VR experiences, and educational AR encounters all share a common temporal structure: visitors engage for seconds to minutes, not hours \cite{vomlehn2001exhibiting, li2024museum}. Ethnographic studies of museum conduct reveal that interactions with exhibits are brief, episodic, and socially organized---visitors negotiate access to objects, point things out to companions, and move on after dwelling for an average of 17--30 seconds \cite{vomlehn2001exhibiting}. Even dedicated exhibition VR installations, where visitors queue for a curated experience, typically offer encounters of one to five minutes before the headset is removed and the next visitor takes a turn.

Yet the devices deployed in these contexts are designed for sustained wearing. Head-mounted displays---whether tethered VR headsets at exhibition stations or AR glasses in museum pilots---are optimized for session-length use: their weight distribution, strap systems, and thermal management assume minutes-to-hours of continuous wear \cite{starner2001challenges}. This creates a practical mismatch. Visitors must navigate complex donning procedures (adjusting interpupillary distance, tightening straps, positioning displays) for encounters that may last only seconds. The ratio of setup time to experience time can approach or exceed 1:1, undermining the spontaneity that makes transient MR encounters compelling.

The mismatch extends into social dynamics. A substantial body of research documents the social costs of wearing head-mounted devices in public. Koelle et al.'s comprehensive survey of social acceptability methods in HCI identifies device visibility, perceived purpose, and wearing context as key determinants of bystander comfort \cite{koelle2020social}. Studies of data glasses adoption show that even extended familiarization fails to fully overcome negative attitudes toward always-worn face-covering devices \cite{koelle2017acceptability}. Profita et al. demonstrate that body location significantly affects the perceived social acceptability of on-body technology, with head and face locations provoking the strongest reactions \cite{profita2013wrist}. Even when HMDs serve clearly beneficial purposes---such as assistive technology for people with visual impairments---social tensions persist, with wearers reporting feelings of conspicuousness and bystanders expressing discomfort \cite{profita2016at, goodman2019social}.

Critically, these social costs are modulated by duration. A two-second peek through a museum viewfinder and a two-hour VR session at a caf\'{e} occupy fundamentally different positions on the social acceptability spectrum. Koelle et al.'s field evaluation of wearable camera acceptability found that bystander comfort depends heavily on perceived recording duration \cite{koelle2019camera}. Kelly and Gilbert's framework identifies the wearer, the device, and its use as determinants of acceptability, but does not formalize duration as a distinct variable \cite{kelly2018wearer}. The social acceptability literature has primarily examined continuous wearing scenarios, leaving the acceptability profile of brief, purposeful wearing largely unexplored.

This creates a paradox: the contexts where MR is most commonly encountered by the public---museums, exhibitions, demonstrations---are precisely the contexts where sustained-wear device assumptions are least appropriate. Transient encounters are forced into devices designed for sessions, producing friction in both the physical and social senses of the word.

\subsection{Temporal Wearing and Wearability}

The wearable computing field has developed sophisticated frameworks for understanding \textit{where} and \textit{how} devices should be worn, but has paid remarkably little attention to \textit{for how long}. Gemperle et al.'s foundational work on wearability at ISWC '98 proposed thirteen design guidelines organized around the spatial relationship between body and device: placement, form language, movement accommodation, sizing, attachment, weight distribution, and aesthetics \cite{gemperle1998wearability}. These guidelines have shaped two decades of wearable design, yet they implicitly assume extended wearing. The thirteenth guideline, ``long-term use,'' addresses durability and hygiene for prolonged wear; no corresponding guideline addresses the design challenges of \textit{brief} or \textit{transitional} wearing.

Zeagler's body maps for wearable placement extended Gemperle's framework with functional, technical, and social considerations for each body location \cite{zeagler2017where}. Social factors touch on duration indirectly---a wrist-worn device is more socially acceptable partly because wrist interactions are brief and familiar---but duration itself is not formalized as a design dimension. Starner's influential characterization of wearable computing challenges distinguishes always-available from task-specific devices \cite{starner2001challenges}, coming closest to a temporal distinction, but task-specific still implies session-length engagement of minutes to hours. Knight et al.'s comfort assessment framework includes ``ease of donning and doffing'' as a usability factor \cite{knight2006comfort}, but treats it as friction to be minimized---an overhead cost---rather than a design opportunity.

The closest precedent to temporal wearing as a design concept comes from Khurana et al.'s work on the detachable smartwatch \cite{khurana2019detachable}. Their IMWUT paper explores a smartwatch that can transition between wrist-worn and handheld modes, demonstrating that the same device can operate in different wearing modalities. This implicitly acknowledges that wearing relationships are not fixed but dynamic. However, Khurana et al. focus on device versatility---a single device serving multiple form factors---rather than on wearing duration as a design dimension. The temporal dimension of when and for how long the device is in each mode remains unexplored.

What makes this gap particularly striking is that temporal wearing has deep roots in optical culture. Long before wearable computing, humans designed and used optical devices intended for seconds-long encounters. The lorgnette, popular from the 18th through early 20th centuries, consisted of a pair of lenses mounted on a handle. Users raised it to their eyes for a focused look---at a companion across a ballroom, at a painting in a gallery, at an actor on stage---and lowered it. The gesture of raising and lowering was socially codified, even ritualized. Opera glasses served an analogous function in theatrical contexts. The monocle, the hand-held magnifying glass, the stereoscope and its parlor-room cousin the View-Master, coin-operated binoculars at scenic overlooks, and the camera viewfinder: all are temporal wearing devices. They were designed not for sustained use but for a rhythm of engagement---raise, look, lower, move, raise again.

These historical form factors embody design principles that contemporary wearable computing has largely forgotten. They optimize for instant-on activation (no boot sequence, no calibration), rapid body coupling (raise to eyes in under one second), social legibility (the gesture of looking through a device is universally understood), and graceful doffing (lowering the device is as natural as raising it). They treat the don/doff cycle not as a cost but as the fundamental unit of interaction. A lorgnette user does not ``set up'' and ``tear down'' a viewing session; each raise-and-lower \textit{is} the session.

Park et al.'s work on comfort and user acceptance establishes that acceptable weight, volume, and contact area thresholds vary by body location \cite{park2019comfort}. Viewed through a temporal lens, these thresholds also vary by \textit{duration}: users tolerate substantially more weight and bulk for a three-second peek than for a three-hour session. This suggests that temporal wearing relaxes certain comfort constraints while introducing others---particularly the need for instant activation and information density within the first second of wear.

The gap, then, is this: wearability research has extensively theorized the spatial dimension of wearing (where on the body, how attached, what form factor) but has not theorized the temporal dimension (for how long, how often, and crucially, whether the wearing duration itself constitutes the interaction). Temporal wearing fills this gap by proposing duration as a primary design variable and identifying a category---seconds-long wearing where don/doff is the interaction---that existing taxonomies do not accommodate.

\subsection{Transitional Interaction as a Cyclic Process}

Temporal wearing is fundamentally about transitions: transitions between not-wearing and wearing, between unaugmented and augmented perception, between individual and shared experience. Several research traditions have examined transitional phenomena in HCI, but none has captured the cyclic, repeating character of temporal wearing.

Weiser and Brown's vision of calm technology proposed that the most effective technologies ``engage both the center and the periphery of our attention, and in fact move back and forth between the two'' \cite{weiser1996calm}. This center-periphery model operates at the information level: a stock ticker in peripheral vision, a notification that briefly demands focal attention. Temporal wearing introduces a \textit{device-level} analog to this information-level transition. The device itself moves between absence (resting on a surface, hanging from a neck strap) and presence (raised to the eyes, held against the body). It is calm in Weiser and Brown's sense---it demands attention only briefly and purposefully, then recedes---but the transition is more dramatic than they envisioned: not a shift of attention within a persistent information environment, but a physical engagement and disengagement with a perceptual device.

Bakker et al.'s work on peripheral interaction formalizes this attention dimension \cite{bakker2015peripheral}. They characterize peripheral interaction as occurring in the ``background or periphery of attention,'' with low attention demand and seamless integration with primary activities. Temporal wearing creates a complementary pattern: a \textit{deliberate, bounded transition} from peripheral to focal attention and back. Unlike peripheral interaction, which aims to keep information in the periphery, temporal wearing momentarily foregrounds the device and its augmented content, then returns to a peripheral or absent state. The result is a punctual attention shift rather than a sustained one---what we might call a \textit{perceptual punctuation mark} in the flow of experience.

The concept of micro-mobility, introduced by Luff and Heath in the context of collaborative work, describes the fine-grained, moment-by-moment repositioning of artifacts---documents, objects, tools---within a shared workspace \cite{luff1998mobility}. People tilt, rotate, push, and pull objects to facilitate shared viewing and collaborative action. Heath et al. extend this to technology-mediated settings, examining how devices feature in the contingent, situated conduct of everyday interaction \cite{heath2000technology}. Temporal wearing extends micro-mobility from documents and artifacts to \textit{devices and the body}: the act of raising a neck-hung viewer to the eyes, tilting it for a companion to see, passing a handheld device to another person, and setting it back on a surface are all micro-mobility actions with a temporal wearing device. The device's position relative to the body is not fixed but continuously negotiated through small physical adjustments.

Recent work on transitional and hybrid interfaces has examined how users move between devices, modalities, and interaction states. Hubenschmid et al.'s survey of hybrid user interfaces in MR catalogs patterns of complementary use across 2D devices and MR environments, identifying sequential, parallel, and complementary transition types \cite{hubenschmid2025hybrid}. Ens et al.'s work on candid interaction addresses the visibility and hiddenness of wearable computing activities, proposing techniques for managing the social legibility of device use \cite{ens2015candid}. Dhaka et al. examine AR user interface transitions at ISMAR, focusing on how AR interfaces transition between visual states \cite{dhaka2024transitions}. These contributions address important aspects of transitional interaction, but they share a common assumption: transitions are \textit{linear}---from state A to state B, from device 1 to device 2, from hidden to visible. They do not account for transitions that are \textit{cyclic}---from not-wearing to wearing to not-wearing to wearing again---where the cycle itself is the fundamental unit of the experience.

Suchman's foundational work on situated action provides a theoretical grounding for why temporal wearing interactions resist linear modeling \cite{suchman2007human}. In Suchman's framework, action is not the execution of pre-formed plans but emerges from the contingent interaction between a person and their context. Temporal wearing is paradigmatically situated: the decision to pick up a device is occasioned by noticing an interesting exhibit, being prompted by a companion, or seeing another visitor's reaction. Each wearing episode is a fresh response to contextual circumstances, not a step in a predetermined sequence. This situatedness distinguishes temporal wearing from habitual wearing (permanent wear, where the device is donned once and forgotten) and planned wearing (session wear, where the user commits to an extended experience in advance).

\subsubsection{The Gap: Cyclic Transition in Temporal Wearing.}
The transition literature reveals a consistent pattern: transitions are modeled as movements between two states---center and periphery \cite{weiser1996calm}, focal and peripheral attention \cite{bakker2015peripheral}, one device and another \cite{hubenschmid2025hybrid}, visible and hidden \cite{ens2015candid}. What is missing is a model of transition as a \textit{cycle}: a repeating rhythm of engagement and disengagement where each cycle is complete in itself yet invites the next. Temporal wearing introduces precisely this cyclic structure. The interaction is not attract $\rightarrow$ don $\rightarrow$ use $\rightarrow$ doff (a linear sequence) but attract $\rightarrow$ don $\rightarrow$ glimpse $\rightarrow$ doff $\rightarrow$ rest $\rightarrow$ return (a cycle that may repeat multiple times). The ``rest'' and ``return'' phases---where the user sets the device down but remains nearby, processes what they saw, and decides whether to engage again---are as integral to the interaction as the glimpse itself.

\subsection{Summary: Three Gaps, One Concept}

Three observations emerge from this review, each supported by extensive prior work yet never connected:

\begin{enumerate}
    \item \textbf{MR is increasingly transient.} The most common public encounters with mixed reality---in museums, exhibitions, demonstrations, and educational settings---are measured in seconds, yet the devices serving these encounters are designed for sustained sessions \cite{vomlehn2001exhibiting, koelle2020social, starner2001challenges}.
    \item \textbf{Optical wearing has always been temporal.} Centuries of optical device design---lorgnettes, opera glasses, stereoscopes, viewfinders---embody a seconds-long wearing pattern, but wearable computing's spatial frameworks have not theorized this temporal dimension \cite{gemperle1998wearability, zeagler2017where}.
    \item \textbf{The interaction is cyclic, not linear.} Users engage with temporal wearing devices through repeating cycles of donning and doffing, but existing transition frameworks model only linear state changes \cite{weiser1996calm, bakker2015peripheral, hubenschmid2025hybrid}.
\end{enumerate}

\textit{Temporal wearing} unifies these three observations into a single concept. It names a widespread but untheorized practice, fills a gap in wearability taxonomies, and introduces a cyclic model of transitional interaction. In the following section, we formalize this concept, propose a wearing duration taxonomy, and define a design space for temporal wearing devices.

% Section 3: Temporal Wearing: Concept & Design Space
% Paper: "Just a Glimpse: Temporal Wearing as a Design Paradigm for Transient Mixed Reality"
% Target: IMWUT May 1, 2026

\section{Temporal Wearing: Concept \& Design Space}

In this section, we introduce \textit{temporal wearing} as a distinct interaction paradigm, present a taxonomy of wearing duration, and articulate the design space that governs the creation of temporally worn mixed reality devices. We ground this framework in both contemporary HCI theory and centuries of precedent in optical instrument design.

\subsection{Defining Temporal Wearing}

We define \textbf{temporal wearing} as an interaction paradigm in which the act of donning, using, and doffing a body-coupled device constitutes the interaction itself, with the entire encounter lasting seconds rather than minutes or hours. In temporal wearing, \textit{wearing duration is the primary design variable}: the device is not a persistent tool that mediates ongoing activity, but a transient lens through which a user briefly accesses an augmented experience before returning to unmediated perception.

This definition rests on a critical distinction from \textit{micro-interactions}~\cite{ashbrook2010microinteractions}. Micro-interactions describe brief input gestures (under four seconds) performed on permanently worn or carried devices---a glance at a smartwatch, a tap on a fitness tracker. The device remains on the body before and after the gesture; the interaction is brief, but the \textit{wearing} is not. Temporal wearing inverts this relationship: the \textit{device encounter itself} is brief. The user was not wearing the device moments ago, will not be wearing it moments hence, and the value of the interaction is inseparable from the physical act of bringing the device to the body and releasing it again.

This inversion has profound design consequences. When donning and doffing are overhead---friction to be endured before productive use begins---designers optimize for sustained comfort, extended battery life, and deep content experiences that justify the transition cost. When donning and doffing \textit{are} the interaction, designers must instead optimize for instant onset, high information density per second, graceful exit, and social legibility. The temporal wearing paradigm thus demands a fundamentally different design logic from both permanent and session-based wearable computing.

Formally, we characterize temporal wearing by three properties:
\begin{enumerate}
    \item \textbf{Brief duration}: The complete don--use--doff cycle lasts on the order of seconds (typically 2--15\,s), not minutes or hours.
    \item \textbf{Donning as interaction}: The physical act of bringing the device to the body is not preparatory overhead but an integral, productive part of the experience.
    \item \textbf{Cyclicity}: Temporal wearing naturally invites repeated encounters---multiple don--doff cycles within a single visit or session---rather than a single extended wearing period.
\end{enumerate}


\subsection{Wearing Duration Taxonomy}

Despite over 25 years of wearable computing research, no prior work has proposed a taxonomy in which \textit{wearing duration} serves as a primary organizing dimension. Gemperle et al.'s foundational framework~\cite{gemperle1998wearability} addresses \textit{where} and \textit{how} to wear devices through 13 design guidelines, but implicitly assumes extended or permanent wear (Guideline~13, ``long-term use''). Zeagler~\cite{zeagler2017where} extends this with body maps for wearable placement, again with an implicit assumption of sustained wearing. Starner~\cite{starner2001challenges} distinguishes always-available from task-specific wearables, but task-specific use still implies session-length engagement. Knight et al.~\cite{knight2006comfort} include ``ease of donning and doffing'' as a comfort metric, treating it as friction to minimize rather than a design opportunity to exploit.

We propose a three-part taxonomy that foregrounds wearing duration as a design dimension (Table~\ref{tab:duration-taxonomy}):

\begin{table*}[t]
\centering
\caption{A taxonomy of wearing duration. Temporal wearing occupies a distinct design region with qualitatively different constraints and opportunities from permanent and session wear.}
\label{tab:duration-taxonomy}
\begin{tabular}{p{2.2cm} p{1.8cm} p{3.2cm} p{3.5cm} p{3.8cm}}
\toprule
\textbf{Category} & \textbf{Duration} & \textbf{Don/Doff Relationship} & \textbf{Design Priority} & \textbf{Examples} \\
\midrule
Permanent wear & Hours--days & Don once, wear continuously & Comfort, aesthetics, battery life, unobtrusiveness & Smartwatch, hearing aid, smart glasses, smart ring \\
\addlinespace
Session wear & Minutes--hours & Don/doff as session bookends & Immersion, sustained comfort, content depth & VR headset, AR HMD, surgical loupes, sports tracker \\
\addlinespace
\textbf{Temporal wear} & \textbf{Seconds} & \textbf{Don/doff IS the interaction} & \textbf{Instant-on, information density, social legibility} & \textbf{Peek viewer, lorgnette, handheld AR, coin-op viewfinder} \\
\bottomrule
\end{tabular}
\end{table*}

Several properties distinguish temporal wearing from shorter instances of session wear. First, \textit{comfort constraints are relaxed}: users tolerate weight, bulk, and ergonomic imperfection for seconds in ways they would not for minutes~\cite{park2019comfort}. Second, \textit{social legibility is enhanced}: brief, purposeful acts (raising binoculars, peering through a viewfinder) are easier for bystanders to interpret than sustained device use, which may appear asocial or surveillance-like~\cite{koelle2020social}. Third, \textit{context triggers wearing}: temporal wearing is \textit{occasioned}---a situated response to environmental stimuli~\cite{suchman2007humanmachine}---rather than habitual. Finally, \textit{design must maximize information per second}: there is no time for tutorials, onboarding, calibration, or gradual content revelation.


\subsection{Design Space Dimensions}

We articulate the temporal wearing design space along three orthogonal dimensions: \textit{body coupling}, \textit{transition cost}, and \textit{social legibility}.

\subsubsection{Body Coupling.}
Body coupling describes the spatial relationship between the user's body and the device during use, ranging from loose to tight coupling:
\begin{itemize}
    \item \textbf{Handheld}: The device is grasped and brought to the eyes or face (e.g., a lorgnette, handheld viewer, or phone-based AR). Coupling is volitional and instantly reversible.
    \item \textbf{Neck-hung}: The device rests on the body passively (e.g., hung from a lanyard) and is raised to the eyes when needed. Coupling alternates between passive carry and active use.
    \item \textbf{Head-mounted}: The device is placed on the head (e.g., a headset or mask). Coupling requires deliberate donning and creates a more committed relationship.
    \item \textbf{Stationary}: The device is fixed in space, and the user approaches with their body (e.g., a coin-operated viewfinder or a mounted stereoscope). The body couples to the device rather than vice versa.
\end{itemize}

Looser body coupling generally enables faster transitions and lower commitment, making it more naturally compatible with temporal wearing. However, tighter coupling can provide richer perceptual experiences (wider field of view, binocular depth, spatial audio) that may justify slightly longer temporal encounters.

\subsubsection{Transition Cost.}
Transition cost quantifies the temporal and physical effort required to move between the unworn and worn states. We operationalize it as:
\begin{equation}
    T_{\text{transition}} = T_{\text{onset}} + T_{\text{offset}}
\end{equation}
where $T_{\text{onset}}$ is the time from initiating donning to receiving the first meaningful visual frame, and $T_{\text{offset}}$ is the time from initiating doffing to fully returning to unmediated perception. For temporal wearing to function as a paradigm, transition cost must be low relative to use duration---ideally under two seconds total, ensuring that at least half the encounter involves augmented content rather than mechanical overhead.

Transition cost is influenced by body coupling (handheld devices have near-zero onset), device design (optical alignment, display wake time, tracking initialization), and physical form factor (straps, clasps, weight distribution). Crucially, transition cost is not merely a usability metric to be minimized: it defines the \textit{boundary of possibility} for temporal wearing. Devices with transition costs exceeding 10--15 seconds effectively preclude temporal use, as the overhead eclipses any brief content experience.

\subsubsection{Social Legibility.}
Social legibility describes how clearly bystanders can perceive and interpret the user's interaction with the device. We decompose it into three components:
\begin{itemize}
    \item \textbf{Action legibility}: Can bystanders see what the user is physically doing? (Raising binoculars is legible; tapping a hidden earpiece is not.)
    \item \textbf{Purpose legibility}: Can bystanders infer \textit{why} the user is performing the action? (Looking through a viewfinder at a scenic overlook is purpose-legible; wearing opaque goggles in a corridor is not.)
    \item \textbf{Duration legibility}: Can bystanders anticipate that the interaction will be brief? (A quick peek through a handheld device signals transience; strapping on a headset does not.)
\end{itemize}

Social legibility is critical for temporal wearing in public and shared spaces. The social acceptability literature establishes that device visibility, perceived purpose, and wearing duration all modulate bystander comfort~\cite{koelle2020social, profita2013wrist}. Temporal wearing, by its nature, optimizes duration legibility: the brevity of the encounter itself communicates that the user has not ``checked out'' of social copresence. Designers can further enhance legibility through transparent or semi-transparent optics, visible content (screens rather than occluded optics), and form factors that reference familiar objects (binoculars, cameras, magnifying glasses).


\subsection{Design Space Mapping}

Figure~\ref{fig:design-space} maps our six prototypes alongside existing commercial and research devices on a two-dimensional plot of \textit{body coupling} ($x$-axis, from handheld to stationary) versus \textit{transition cost} ($y$-axis, from near-zero to 30+ seconds). Each device is annotated with its social legibility (represented by marker size, where larger markers indicate higher legibility).

\begin{figure}[t]
    \centering
    % Placeholder for design space figure
    % 2D scatter plot: x = body coupling (handheld -> neck-hung -> head-mounted -> stationary)
    % y = transition cost (ms, log scale or linear 0-30s)
    % Marker size = social legibility
    % Our 6 prototypes: cluster in low-transition-cost region across coupling types
    % Reference devices: Meta Quest (head-mounted, high transition), HoloLens (head-mounted, high),
    %   Google Glass (head-mounted but lower transition), phone AR (handheld, low transition),
    %   coin-op viewfinder (stationary, low transition)
    \includegraphics[width=\columnwidth]{figures/design-space-placeholder}
    \caption{Design space mapping of temporal wearing devices. Our six prototypes (filled markers) cluster in the low-transition-cost region across body coupling types. Existing devices (open markers) are plotted for comparison. Marker size encodes social legibility. The shaded region indicates the \textit{temporal wearing zone}: combinations of body coupling and transition cost that support don--use--doff cycles of under 15 seconds.}
    \label{fig:design-space}
\end{figure}

The mapping reveals that our prototypes occupy a distinct region of the design space---the \textit{temporal wearing zone}---characterized by transition costs under 3 seconds across multiple body coupling modalities. Existing MR devices (Meta Quest, HoloLens~2, Magic Leap~2) cluster in the high-transition-cost, head-mounted quadrant, designed for session wear. Phone-based AR (handheld, low transition cost) falls within the temporal wearing zone but lacks the perceptual immersion of optical see-through designs. Coin-operated viewfinders and stereoscopic viewers (stationary, near-zero transition cost) have long occupied this space without being recognized as temporal wearing devices.


\subsection{The Glimpse Cycle}

Session-based wearable interaction follows a linear trajectory: \textit{don $\rightarrow$ use $\rightarrow$ doff}. The user commits to wearing the device, engages in extended activity, and eventually removes it. Temporal wearing, by contrast, follows a \textit{cyclic} pattern that we term the \textbf{glimpse cycle} (Figure~\ref{fig:glimpse-cycle}):

\begin{figure}[t]
    \centering
    % Placeholder for glimpse cycle figure
    % Circular diagram: attract -> approach -> don -> glimpse -> doff -> (rest) -> return -> attract...
    \includegraphics[width=\columnwidth]{figures/glimpse-cycle-placeholder}
    \caption{The glimpse cycle. Unlike the linear don--use--doff model of session wear, temporal wearing follows a repeating cycle in which each phase carries distinct design implications.}
    \label{fig:glimpse-cycle}
\end{figure}

\begin{enumerate}
    \item \textbf{Attract}: The user notices the device or its effects---a visual cue, another user's reaction, signage, or the device's physical presence. Design implication: the device must be \textit{discoverable} and \textit{inviting}.
    \item \textbf{Approach}: The user moves toward the device (or picks it up). Design implication: the device's affordances must communicate how it is used from a distance.
    \item \textbf{Don}: The user brings the device into contact with the body (raises it to their eyes, places it on their head, leans into a viewfinder). Design implication: the onset must be \textit{instant}---content should appear within one second of physical contact.
    \item \textbf{Glimpse}: The user experiences augmented content. This is the core perceptual moment, lasting 2--10 seconds. Design implication: content must deliver value \textit{immediately}, without onboarding or orientation.
    \item \textbf{Doff}: The user lowers or removes the device. Design implication: exit must feel \textit{natural and graceful}, not abrupt or disorienting.
    \item \textbf{Rest}: The user returns to unmediated perception. They may process what they saw, share it socially, or simply continue their activity. Design implication: the augmented experience should leave a \textit{residue}---a prompt for reflection, conversation, or return.
    \item \textbf{Return}: The user re-engages with the device for another glimpse. Design implication: content should \textit{reward re-engagement} (e.g., through spatial variation, temporal change, or progressive revelation) rather than penalizing brevity.
\end{enumerate}

This cyclicity is a defining feature of temporal wearing. Where session wear typically involves a single don--use--doff cycle per engagement, temporal wearing invites multiple cycles within a single visit. Our field observations (Section~5) confirm that users of temporally worn devices frequently engage in 3--8 glimpse cycles per visit, with rest periods ranging from seconds to minutes. This pattern more closely resembles how people use binoculars at a scenic overlook or flip through a View-Master reel than how they use a VR headset.

The glimpse cycle also has social implications. The rest phase creates windows of social availability---moments when the user is visibly present and interactable. This contrasts sharply with session wear, where the extended don phase creates prolonged social absence. Temporal wearing thus supports what we call \textit{shared glimpsing}: a social mode in which multiple users take turns engaging with a device, narrate their experiences to companions, and collectively construct meaning from individual glimpses.


\subsection{Historical Precedents}

Temporal wearing is not a novel invention but a rediscovery. Optical instruments designed for brief, repeated viewing---raising to the eyes, glimpsing, and lowering---have existed for centuries. These precedents demonstrate that temporal wearing is a natural and culturally established interaction pattern, long practiced but never formally named within HCI.

\textbf{The lorgnette} (circa 1770s--1920s) was a pair of spectacles held on a handle rather than worn on the ears~\cite{rosenthal1996spectacles}. Users raised the lorgnette to their eyes momentarily to read, examine, or observe, then lowered it---a paradigmatic temporal wearing interaction. The lorgnette's social function was as important as its optical one: it signaled attention, interest, and social status through the visible act of looking.

\textbf{Opera glasses} (from the 1780s onward) extended the lorgnette's temporal wearing pattern to entertainment contexts~\cite{ilardi2007renaissance}. Theatergoers raised compact binoculars to their eyes for brief periods during performances, alternating between augmented (magnified) and unaugmented viewing. The don--glimpse--doff cycle was integral to the social ritual of theater attendance.

\textbf{The stereoscope} (invented by Charles Wheatstone in 1838, popularized by David Brewster's lenticular design in 1849) introduced binocular depth to temporal viewing~\cite{wade2002biography}. Users held or leaned into the stereoscope to view a three-dimensional image, then withdrew---a peek interaction lasting seconds. The Holmes stereoscope (1861) became a parlor entertainment device, passed among guests for repeated brief viewings: a social glimpse cycle.

\textbf{Coin-operated viewfinders} (introduced in the early 20th century, popularized mid-century) represent stationary temporal wearing: the user approaches a fixed device, leans in, views a magnified or augmented panorama, and steps back. The interaction is inherently brief (often coin-timed to 1--2 minutes, but individual viewing episodes within that window are seconds-long as users pan across a scene). The viewfinder's fixed placement and familiar form make it maximally socially legible.

\textbf{The View-Master} (introduced by William Gruber and Harold Graves in 1939) is perhaps the purest consumer temporal wearing device~\cite{zone20073d}. Users raise the binocular viewer to their eyes, view a stereoscopic image, advance the reel, and lower the device---a rapid glimpse cycle measured in seconds per image. The View-Master's cultural longevity (produced continuously for over 80 years) testifies to the naturalness of temporal wearing as an interaction pattern.

These historical examples share several properties that define the temporal wearing paradigm: (1)~the device is raised to the body and lowered within seconds; (2)~the physical act of looking through the device is inseparable from the content experience; (3)~social context shapes and is shaped by the wearing act; and (4)~repeated brief encounters are the norm, not the exception. Contemporary mixed reality devices, designed overwhelmingly for session wear, have largely abandoned this centuries-old pattern. Temporal wearing, as we propose it, recovers and extends this tradition for the age of spatial computing.

% Section 4: Method
% Paper: "Just a Glimpse: Temporal Wearing as a Design Paradigm for Transient Mixed Reality"

\section{Method}

We adopted a Research through Design (RtD) approach~\cite{zimmerman2007research} to investigate temporal wearing as a design paradigm for mixed reality. RtD is well-suited to this inquiry because our research questions concern not only \textit{whether} temporal wearing produces distinct interaction patterns, but \textit{how} design decisions---body coupling, transition cost, form factor---shape those patterns in situ. Rather than testing a single device against a baseline, we designed a family of prototypes that instantiate different positions in the temporal wearing design space (Section~3.3), deployed a subset in a public setting, and analyzed the resulting interactions through mixed methods. This approach follows the precedent of foundational wearable computing research~\cite{gemperle1998wearability} and recent IMWUT work on novel wearable form factors~\cite{khurana2019detachable}, which establish design principles through iterative prototyping and empirical grounding.

\subsection{Prototypes}

We designed and built six functional MR prototypes spanning three form factors---handheld, neck-hung, and stationary---each optimized for temporal wearing (Table~\ref{tab:prototypes}). All prototypes use optical see-through or video see-through displays to deliver spatially registered augmented content, and all were designed to achieve sub-two-second transition costs (onset + offset). Detailed technical specifications, construction processes, and design rationale for each prototype appear in Appendix~A.

\begin{table*}[t]
\centering
\caption{Overview of six temporal wearing prototypes across three form factors. Transition cost is the measured mean time from initiating donning to first augmented frame (onset) plus mean time from initiating doffing to full return to unmediated perception (offset). Three prototypes (marked with $\dagger$) were selected for the field deployment.}
\label{tab:prototypes}
\begin{tabular}{p{2.8cm} p{1.8cm} p{4.5cm} p{2.2cm} p{2.5cm}}
\toprule
\textbf{Prototype} & \textbf{Form Factor} & \textbf{Key Feature} & \textbf{Transition Cost} & \textbf{Deployment} \\
\midrule
[PROTOTYPE A]$\dagger$ & Handheld & [PLACEHOLDER: e.g., monocular viewer with grip handle] & [PLACEHOLDER] s & Field deployment \\
\addlinespace
[PROTOTYPE B] & Handheld & [PLACEHOLDER: e.g., binocular peek device] & [PLACEHOLDER] s & Lab evaluation \\
\addlinespace
[PROTOTYPE C]$\dagger$ & Neck-hung & [PLACEHOLDER: e.g., lanyard-suspended viewer] & [PLACEHOLDER] s & Field deployment \\
\addlinespace
[PROTOTYPE D] & Neck-hung & [PLACEHOLDER: e.g., pendant-style AR monocle] & [PLACEHOLDER] s & Lab evaluation \\
\addlinespace
[PROTOTYPE E]$\dagger$ & Stationary & [PLACEHOLDER: e.g., plinth-mounted viewfinder] & [PLACEHOLDER] s & Field deployment \\
\addlinespace
[PROTOTYPE F] & Stationary & [PLACEHOLDER: e.g., wall-mounted stereoscope] & [PLACEHOLDER] s & Lab evaluation \\
\bottomrule
\end{tabular}
\end{table*}

The prototypes were designed to systematically vary body coupling while holding content constant: all six present the same augmented content experience, a [PLACEHOLDER: brief description of AR content, e.g., spatially registered overlay revealing hidden layers of an artwork]. This controlled variation allows us to isolate the effects of form factor and transition cost on temporal wearing behavior. Each form factor pair includes one prototype optimized for minimal transition cost and one that explores a different point in the design space (e.g., wider field of view at the cost of slightly longer onset time).

\subsection{Field Deployment Design}

From the six prototypes, we selected three for field deployment---one per form factor---to investigate the following research questions:

\begin{itemize}
    \item[\textbf{RQ1}] How does transition cost affect re-engagement frequency? We hypothesize that prototypes with lower transition costs will invite more frequent glimpse cycles within a single visit, as the reduced overhead makes repeated donning and doffing more natural.
    \item[\textbf{RQ2}] How does body coupling affect social comfort in co-located settings? We hypothesize that looser body coupling (handheld, stationary) will produce higher social comfort scores among both users and bystanders, as these form factors maintain social legibility and minimize the appearance of isolation.
    \item[\textbf{RQ3}] What temporal wearing behaviors emerge in situ? Beyond our design space predictions, we seek to identify the behavioral patterns, social dynamics, and appropriation strategies that arise when people encounter temporal wearing devices in a naturalistic setting.
\end{itemize}

These questions address distinct aspects of temporal wearing: RQ1 targets the mechanical dimension (how physical design shapes interaction frequency), RQ2 targets the social dimension (how form factor mediates social acceptability), and RQ3 targets the emergent dimension (what unanticipated behaviors arise in practice).

\subsection{Deployment Site}

We deployed the three prototypes at [VENUE PLACEHOLDER], a [PLACEHOLDER: type of venue, e.g., contemporary art gallery / science museum / university exhibition space] in [PLACEHOLDER: city, country]. The venue was selected for several properties that make it well-suited to studying temporal wearing:

\begin{itemize}
    \item \textbf{Naturalistic transient encounters}: Visitors arrive voluntarily, move freely through the space, and dwell at individual exhibits for seconds to minutes---matching the temporal wearing paradigm's target context.
    \item \textbf{Social co-presence}: Visitors typically attend in pairs or small groups, enabling observation of shared glimpsing, narration, and hand-off behaviors.
    \item \textbf{Spatial configuration}: The venue provided a [PLACEHOLDER: e.g., single room / corridor / open gallery] of approximately [PLACEHOLDER] m$^2$, allowing all three prototypes to be positioned within sight of one another and within the field of view of our video cameras.
    \item \textbf{Visitor diversity}: The venue attracts a broad demographic, including families, couples, student groups, and solo visitors, providing variation in social context and prior technology experience.
\end{itemize}

Each prototype was placed adjacent to a [PLACEHOLDER: artwork / exhibit / installation], with minimal signage---a small label indicating that visitors could pick up or lean into the device to ``see something hidden.'' This deliberate understatement was designed to test the \textit{attract} phase of the glimpse cycle (Section~3.5): how do visitors discover and approach a temporal wearing device without explicit instruction?

\subsection{Participants}

A total of N=[PLACEHOLDER] visitors interacted with the prototypes during the deployment period of [PLACEHOLDER: duration, e.g., three days / one week]. Participation was naturalistic: visitors encountered the prototypes as part of their normal visit to the venue and were not recruited in advance. We collected demographic information only from the subset of visitors who consented to exit interviews (N=[PLACEHOLDER], described below). Of the interviewed participants, [PLACEHOLDER: age range, gender distribution, prior MR experience summary]. The study was approved by [PLACEHOLDER: ethics board] (protocol \#[PLACEHOLDER]).

\subsection{Deployed Prototypes}

Three prototypes were deployed, one per form factor:

\begin{itemize}
    \item \textbf{Handheld} ([PROTOTYPE A]): [PLACEHOLDER: 1--2 sentence description]. This prototype represents the lowest body coupling and fastest transition cost, analogous to historical lorgnettes and handheld viewers.
    \item \textbf{Neck-hung} ([PROTOTYPE C]): [PLACEHOLDER: 1--2 sentence description]. The lanyard maintains passive body coupling when not in active use, reducing the approach phase and enabling rapid re-engagement.
    \item \textbf{Stationary} ([PROTOTYPE E]): [PLACEHOLDER: 1--2 sentence description]. The fixed position eliminates the need to carry or manage the device, with the user's body approaching the device rather than vice versa.
\end{itemize}

All three prototypes presented the same augmented content to control for content effects. The prototypes were counterbalanced across three positions in the venue, with positions rotated [PLACEHOLDER: daily / every two days] to mitigate location effects.

\subsection{Data Collection}

We employed four complementary data collection methods:

\paragraph{Interaction Logs.}
Each prototype was instrumented to record timestamped don and doff events, operationalized as [PLACEHOLDER: e.g., IMU-based lift detection for handheld, proximity sensor activation for stationary, accelerometer threshold for neck-hung]. Logged events included: timestamp, event type (don/doff), duration of each glimpse episode, and inter-episode interval (time between consecutive doff and don events). These logs provided quantitative measures of transition timing, glimpse duration, re-engagement frequency, and total interaction time per visitor.

\paragraph{Video Observation.}
Two cameras were positioned to capture the deployment area from complementary angles: one overview camera capturing the full spatial context (visitor trajectories, social groupings, approach behaviors) and one close-up camera focused on the prototype interaction zone (donning gestures, facial expressions, device manipulation). Video was recorded continuously during venue opening hours, yielding approximately [PLACEHOLDER] hours of footage. Video data was used for behavioral coding and provided context for interpreting interaction log patterns.

\paragraph{Exit Interviews.}
We conducted semi-structured exit interviews with N=[PLACEHOLDER: 15--20] visitors, selected through purposive sampling to ensure variation in age, group composition, and observed interaction patterns (e.g., single-glimpse visitors, multi-cycle visitors, visitors who engaged in hand-off behaviors). Interviews lasted 10--15 minutes and covered: initial discovery and approach, the experience of donning and doffing, awareness of bystanders during use, decisions to re-engage or move on, and comparisons between prototypes for visitors who tried multiple devices. The full interview protocol appears in Appendix~B.

\paragraph{Social Comfort Survey.}
Immediately following their interaction, a subset of visitors (N=[PLACEHOLDER]) completed a brief survey assessing social comfort during device use and willingness to re-engage. The survey used 7-point Likert scales adapted from established social acceptability instruments~\cite{koelle2020social, rico2010gestures} and included items on: felt self-consciousness, perceived bystander reactions, comfort with the physical act of donning/doffing, and likelihood of using the device again. Bystander comfort was assessed through a parallel survey administered to [PLACEHOLDER: N] companions or nearby visitors who observed but did not use the device. The full survey instrument appears in Appendix~C.

\subsection{Analysis}

We employed a mixed-methods analysis approach combining quantitative interaction metrics with qualitative behavioral and interview analysis.

\paragraph{Interaction Metrics.}
Descriptive statistics (mean, median, standard deviation, range) were computed for onset time, offset time, glimpse duration, number of glimpse cycles per visitor, and inter-episode interval, broken down by prototype. To address RQ1, we computed Spearman rank correlations between transition cost and re-engagement frequency across prototypes. Non-parametric tests (Kruskal-Wallis with Dunn's post-hoc comparisons) were used to compare metrics across form factors, given the expected non-normal distribution of timing data.

\paragraph{Video Coding.}
Video data was coded using a scheme adapted from Brignull and Rogers' framework for public display interaction~\cite{brignull2003enticing}, extended to capture temporal wearing-specific behaviors. Coded categories included: discovery behavior (how visitors first noticed the device), approach trajectory, donning gesture, glimpse behavior (scanning, fixating, searching), doffing gesture, post-doff behavior (social sharing, verbal narration, return), and bystander reactions. The coding scheme was developed iteratively: two researchers independently coded [PLACEHOLDER: 20\%] of the video data, discussed disagreements, refined category definitions, and achieved a Cohen's $\kappa$ of [PLACEHOLDER] before coding the remaining footage independently. The full coding scheme appears in Appendix~D.

\paragraph{Thematic Analysis.}
Interview transcripts were analyzed using reflexive thematic analysis following Braun and Clarke's six-phase approach~\cite{braun2006thematic}: familiarization through repeated reading, systematic coding, theme generation, theme review, theme definition, and report writing. We adopted an inductive-deductive orientation: initial codes were generated from the data without reference to the temporal wearing framework, then organized into candidate themes that were iteratively refined in dialogue with the framework. This approach balances sensitivity to participants' lived experience with the theoretical commitments of Research through Design. Two researchers conducted the analysis collaboratively, meeting regularly to discuss coding decisions and theme development.

% Section 5: Results
% Paper: "Just a Glimpse: Temporal Wearing as a Design Paradigm for Transient Mixed Reality"

\section{Results}

Over the [PLACEHOLDER: duration] deployment, N=[PLACEHOLDER] unique visitors interacted with the three prototypes, generating [PLACEHOLDER] logged interaction events across [PLACEHOLDER] hours of video footage. We present findings organized around our three research questions: transition metrics and re-engagement patterns (RQ1), social comfort and behaviors (RQ2), and emergent temporal wearing behaviors (RQ3).

\subsection{Transition Metrics}

Table~\ref{tab:transition-metrics} reports the measured onset, offset, and total transition costs for each deployed prototype, alongside glimpse duration---the time spent viewing augmented content between don and doff.

\begin{table*}[t]
\centering
\caption{Transition metrics by prototype. Onset = time from initiating don to first augmented frame; Offset = time from initiating doff to full return to unmediated perception; Glimpse duration = time viewing augmented content per episode. All values in seconds; reported as mean (SD).}
\label{tab:transition-metrics}
\begin{tabular}{p{2.5cm} p{2.2cm} p{2.2cm} p{2.5cm} p{2.5cm} p{2.5cm}}
\toprule
\textbf{Prototype} & \textbf{Onset (s)} & \textbf{Offset (s)} & \textbf{Total Transition (s)} & \textbf{Glimpse Duration (s)} & \textbf{Efficiency Ratio} \\
\midrule
Handheld & [PLACEHOLDER] & [PLACEHOLDER] & [PLACEHOLDER] & [PLACEHOLDER] & [PLACEHOLDER] \\
Neck-hung & [PLACEHOLDER] & [PLACEHOLDER] & [PLACEHOLDER] & [PLACEHOLDER] & [PLACEHOLDER] \\
Stationary & [PLACEHOLDER] & [PLACEHOLDER] & [PLACEHOLDER] & [PLACEHOLDER] & [PLACEHOLDER] \\
\bottomrule
\end{tabular}
\end{table*}

The \textit{efficiency ratio}---defined as glimpse duration divided by total interaction time (glimpse duration + transition cost)---captures what proportion of the encounter involves augmented content versus mechanical overhead. [PLACEHOLDER: Report that handheld achieved the highest efficiency ratio, with near-instantaneous onset; stationary had the lowest transition cost overall due to zero approach overhead for the device itself; neck-hung occupied a middle position. Report statistical comparisons.]

Across all prototypes, the mean glimpse duration was [PLACEHOLDER] seconds (SD = [PLACEHOLDER]), confirming that temporal wearing interactions are brief---overwhelmingly under 15 seconds---and validating the temporal wearing category in our duration taxonomy (Section~3.2). Notably, [PLACEHOLDER: report any interesting distribution patterns, e.g., bimodal distribution with peaks at ~3s and ~8s suggesting ``quick peeks'' and ``deliberate looks'' as distinct temporal wearing modes].

\subsection{Re-engagement Patterns}

A central prediction of the temporal wearing framework is that lower transition costs should encourage more frequent re-engagement---more glimpse cycles per visit. Figure~\ref{fig:reengagement} plots the number of glimpse cycles per visitor against prototype form factor.

\begin{figure}[t]
    \centering
    \includegraphics[width=\columnwidth]{figures/reengagement-placeholder}
    \caption{Number of glimpse cycles per visitor by prototype. [PLACEHOLDER: Box plot or violin plot showing distribution of re-engagement counts. Expected pattern: handheld > neck-hung > stationary for re-engagement frequency, though stationary may show a distinct pattern due to fixed location.]}
    \label{fig:reengagement}
\end{figure}

[PLACEHOLDER: Report descriptive statistics for re-engagement. Expected findings:]

Visitors using the handheld prototype completed a mean of [PLACEHOLDER] glimpse cycles per visit (SD = [PLACEHOLDER], range = [PLACEHOLDER]), compared to [PLACEHOLDER] for the neck-hung prototype and [PLACEHOLDER] for the stationary prototype. A Kruskal-Wallis test revealed a significant effect of form factor on re-engagement frequency ($H$([PLACEHOLDER]) = [PLACEHOLDER], $p$ = [PLACEHOLDER]). Post-hoc Dunn's tests indicated that the handheld prototype elicited significantly more re-engagements than [PLACEHOLDER] ($p$ = [PLACEHOLDER]), while [PLACEHOLDER].

The inter-episode interval---the time between putting the device down and picking it up again---also varied by prototype. Mean inter-episode intervals were [PLACEHOLDER] seconds for handheld, [PLACEHOLDER] seconds for neck-hung, and [PLACEHOLDER] seconds for stationary. [PLACEHOLDER: Report whether shorter inter-episode intervals correlated with more total glimpse cycles, and note any interesting patterns such as visitors who walked away and returned later versus those who engaged in rapid successive peeks.]

Spearman rank correlation between measured transition cost and re-engagement frequency across individual visitors yielded $\rho$ = [PLACEHOLDER], $p$ = [PLACEHOLDER], [PLACEHOLDER: supporting/not supporting] the hypothesis that lower transition costs encourage more frequent re-engagement (RQ1).

\subsection{Social Behaviors}

Video coding revealed distinctive social dynamics around temporal wearing interactions. We coded [PLACEHOLDER] interaction episodes from [PLACEHOLDER] hours of video, identifying [PLACEHOLDER] episodes involving social interaction (defined as verbal or gestural communication between the user and at least one co-present person during or immediately after device use).

\paragraph{Shared Viewing.}
[PLACEHOLDER: \% of interactions] involved some form of shared viewing---users who were accompanied by at least one other person. Of these, [PLACEHOLDER: \%] involved the companion watching the user interact, [PLACEHOLDER: \%] involved the companion subsequently using the same device, and [PLACEHOLDER: \%] involved simultaneous engagement (one person using the device while narrating to the companion). The handheld prototype showed the highest rate of device sharing ([PLACEHOLDER: \%]), likely because the physical act of passing a handheld object is a deeply familiar social gesture.

\paragraph{Bystander Engagement.}
We adapted Brignull and Rogers'~\cite{brignull2003enticing} engagement zones---\textit{peripheral awareness}, \textit{focal awareness}, and \textit{direct interaction}---to analyze how bystanders progressed toward device use. [PLACEHOLDER: Report the honeypot effect~\cite{wouters2016honeypot}---whether seeing one person use a device attracted others. Report typical progression times from peripheral awareness to direct interaction.]

\paragraph{Social Negotiation.}
In [PLACEHOLDER: \%] of multi-person groups, we observed explicit social negotiation around device use: verbal invitations (``You try it''), gestural offers (extending the device toward a companion), and turn-taking patterns. These negotiations were notably brief---typically under [PLACEHOLDER] seconds---reflecting the low-commitment nature of temporal wearing. Unlike session-wear VR experiences, where ``taking a turn'' implies a minutes-long commitment, temporal wearing turns lasted seconds, reducing the social cost of both offering and accepting.

\subsection{Social Comfort}

Survey responses from N=[PLACEHOLDER] users and N=[PLACEHOLDER] bystanders revealed significant differences in social comfort across form factors (Table~\ref{tab:social-comfort}).

\begin{table}[t]
\centering
\caption{Social comfort scores by prototype (7-point Likert scale, higher = more comfortable). Reported as mean (SD).}
\label{tab:social-comfort}
\begin{tabular}{p{2cm} p{2cm} p{2cm} p{2cm}}
\toprule
\textbf{Measure} & \textbf{Handheld} & \textbf{Neck-hung} & \textbf{Stationary} \\
\midrule
User self-consciousness & [PH] & [PH] & [PH] \\
User social comfort & [PH] & [PH] & [PH] \\
Bystander comfort & [PH] & [PH] & [PH] \\
Re-engagement willingness & [PH] & [PH] & [PH] \\
\bottomrule
\end{tabular}
\end{table}

[PLACEHOLDER: Report statistical comparisons. Expected pattern: handheld and stationary score highest on social comfort (familiar gestures---picking up an object, leaning into a viewfinder); neck-hung may score slightly lower due to unfamiliarity of the form factor. Report any interesting interactions between social comfort and re-engagement willingness.]

Qualitatively, several survey respondents commented on the relationship between brevity and comfort: [PLACEHOLDER: representative quotes, e.g., ``It didn't feel weird because it was so quick'' or ``I would have felt self-conscious wearing a headset, but picking this up felt natural''].

\subsection{Emergent Temporal Wearing Behaviors}

Beyond the predicted effects of transition cost and body coupling, our video analysis and interviews revealed four distinctive temporal wearing behaviors that emerged spontaneously across the deployment (RQ3). We name these behaviors for their characteristic patterns.

\subsubsection{``The Double-Take.''}
In [PLACEHOLDER: N / \%] of interactions, visitors performed what we call a \textit{double-take}: a very brief initial glimpse (under [PLACEHOLDER] seconds), followed by lowering the device, a pause of [PLACEHOLDER] seconds, and then a second, longer glimpse. Interview data suggests this pattern reflects a two-stage cognitive process: the first glimpse establishes \textit{that} augmented content exists (a perceptual surprise), and the second glimpse explores \textit{what} the content contains (directed attention). As participant P[PLACEHOLDER] described:

\begin{quote}
``The first time I just went---oh! There's something there. And then I had to look again properly to actually see what it was.'' (P[PLACEHOLDER])
\end{quote}

The double-take was most common with the handheld prototype ([PLACEHOLDER: \%]), where the near-zero transition cost made the ``peek and return'' gesture effortless. It was less common with the stationary prototype ([PLACEHOLDER: \%]), where the physical commitment of leaning in may have encouraged longer initial glimpses.

\subsubsection{``The Hand-Off.''}
In [PLACEHOLDER: N / \%] of group interactions, one visitor used a prototype and then immediately handed it (or directed their companion toward it) to another person. We term this \textit{the hand-off}---a social sharing behavior in which one person's glimpse directly occasions another's. The hand-off was almost exclusively observed with the handheld prototype ([PLACEHOLDER: \%] of handheld group interactions), where the physical act of passing the device is natural and the gesture itself communicates ``you should see this.'' The neck-hung prototype produced a variant: rather than passing the device, the wearer would lift it from their own chest and hold it toward the companion's face---a hybrid hand-off that maintained the lanyard connection.

\begin{quote}
``I saw her reaction and I just had to see what she was seeing. She handed it to me without me even asking.'' (P[PLACEHOLDER])
\end{quote}

The hand-off behavior connects temporal wearing to the social dynamics of museum visiting documented by vom Lehn et al.~\cite{vomlehn2001exhibiting}: visitors routinely draw companions' attention to exhibits through pointing, verbal direction, and physical repositioning. Temporal wearing devices, by being hand-offable, integrate into this existing social repertoire.

\subsubsection{``The Narrator.''}
In [PLACEHOLDER: N / \%] of group interactions, we observed a distinctive division of labor: one person wore or held the device and simultaneously described what they saw to a companion who was not using any device. We call this \textit{the narrator} pattern. Unlike the hand-off, where the device moves between people, the narrator pattern creates a shared experience through verbal and gestural narration: the user describes what they see (``Oh, the colors are moving---there's a blue light coming from the painting''), and the companion responds, asks questions, and participates in meaning-making without ever donning the device.

\begin{quote}
``I didn't need to look myself. She was telling me everything, and it was actually more fun watching her face and hearing her describe it.'' (P[PLACEHOLDER])
\end{quote}

The narrator pattern is particularly significant because it demonstrates that temporal wearing can produce \textit{socially shared MR experiences without requiring multiple devices or simultaneous wearing}. The brevity of temporal wearing---and the fact that the user remains visibly present and conversationally available---makes narration possible in ways that sustained HMD use does not. A person wearing a VR headset for five minutes is socially absent; a person peeking through a handheld viewer for five seconds while describing what they see is socially \textit{present}.

\subsubsection{``The Orbit.''}
In [PLACEHOLDER: N / \%] of all interactions, visitors returned to a prototype after having moved on to other parts of the venue. We call this \textit{the orbit}---a large-scale re-engagement pattern in which visitors circled back to a device during their visit, sometimes multiple times. The orbit was observed across all form factors but was most frequent with the stationary prototype ([PLACEHOLDER: \%]), perhaps because its fixed location made it a spatial landmark that visitors could easily find again.

Orbit intervals ranged from [PLACEHOLDER] to [PLACEHOLDER] minutes, with some visitors returning two or three times over the course of a [PLACEHOLDER]-minute visit. Interview data suggests multiple motivations for orbiting: wanting to see if the content had changed (``I wondered if it would be different the second time''), wanting to show a new companion (``My friend arrived later and I wanted her to see it''), and simply wanting to re-experience the perceptual pleasure of the glimpse (``It was just nice to look'').

\begin{quote}
``I kept coming back. Each time I noticed something different. It was like the painting kept revealing itself.'' (P[PLACEHOLDER])
\end{quote}

\subsection{Interview Themes}

Reflexive thematic analysis of exit interviews (N=[PLACEHOLDER]) generated [PLACEHOLDER] themes organized around three overarching categories: the experience of temporal wearing, social dimensions of device use, and comparisons with prior MR encounters.

\subsubsection{Theme 1: The Pleasure of Brevity.}
Participants consistently described the brief encounter as a positive quality rather than a limitation. Unlike session-based MR experiences, where users may feel pressure to ``get their money's worth'' or ``explore everything,'' temporal wearing encounters were described as \textit{gifts}---small, self-contained perceptual pleasures with no obligation to continue. [PLACEHOLDER: 2--3 supporting quotes.]

\subsubsection{Theme 2: Wearing Without Commitment.}
Several participants contrasted their temporal wearing experience with prior HMD use, emphasizing the low commitment involved: ``I didn't have to \textit{become} someone wearing a headset'' (P[PLACEHOLDER]). This theme captures a felt distinction between \textit{being a device wearer} (an identity, however temporary) and \textit{using a device momentarily} (an action). Temporal wearing sits firmly in the latter category, which participants associated with reduced self-consciousness and greater willingness to engage.

\subsubsection{Theme 3: Social Permission.}
Participants in groups described a social dynamic in which one person's use of the device gave ``permission'' for others to try it. This \textit{social permission} effect was amplified by the visibility and brevity of temporal wearing: seeing someone pick up a device, peek through it for a few seconds, and put it down made the interaction appear safe, simple, and brief---lowering the threshold for companions to try. [PLACEHOLDER: supporting quotes.]

\subsubsection{Theme 4: The Hidden World.}
Multiple participants described a sense of \textit{discovery}---that the augmented content was a ``hidden world'' or ``secret layer'' only visible through the device. This discovery framing was enhanced by the temporal nature of the encounter: the augmented world was not persistent but \textit{glimpsed}, making each viewing feel like an act of revelation rather than sustained observation. [PLACEHOLDER: supporting quotes.]

\subsubsection{Theme 5: Wanting More (in the Right Way).}
Participants expressed a desire to return and see more, but framed this as anticipatory pleasure rather than frustration: ``I want to come back and look again'' rather than ``I wish I could have looked longer.'' This finding suggests that temporal wearing, when designed well, creates an \textit{invitation to return} rather than a sense of deprivation---a crucial distinction for the design principles we derive in Section~6. [PLACEHOLDER: supporting quotes.]

% Section 6: Discussion
% Paper: "Just a Glimpse: Temporal Wearing as a Design Dimension for Transient Mixed Reality"

\section{Discussion}

Our findings across both study phases suggest that temporal wearing is not merely a shorter version of session wear but a qualitatively different mode of interaction with its own behavioral patterns, social dynamics, and design requirements. In this section, we reflect on wearing duration as a design dimension, discuss what the baseline comparison revealed, articulate six design principles derived from our empirical findings, consider applications beyond mixed reality, and acknowledge limitations.

\subsection{Wearing Duration as a Design Dimension}

The dominant trajectory of MR development treats progress as longer, more immersive, more seamless wearing. Each generation of HMDs is lighter, more comfortable, and more capable---all in service of extending wearing duration. Apple Vision Pro's ``all-day'' comfort, Meta Quest 3's extended battery life, and the industry's pursuit of lightweight AR glasses all reflect this logic: the ideal device is one you never take off. Our findings suggest a complementary perspective---not that sustained wear is wrong, but that wearing duration is a consequential design dimension that has been undertheorized, and that the brief end of the duration spectrum has distinct properties worth designing for.

The lab study (Phase~1) provides initial evidence for this claim. When participants used the same content through temporal wearing versus session wear, the two modes appeared to produce different interaction patterns: temporal wearing conditions elicited more frequent re-engagement cycles, and participants reported [PLACEHOLDER: higher/comparable] social comfort with temporal wearing forms. The deployment (Phase~2) extends this finding into a naturalistic setting, where the emergent behaviors---double-takes, hand-offs, narration, orbiting---suggest that temporal wearing invites interaction patterns that session wear does not readily afford.

When participants performed double-takes, hand-offs, narrations, and orbits, they were not ``failing to sustain engagement''---they were engaging in precisely the way temporal wearing invites. The efficiency ratio data (Section~5.1) indicates that temporal wearing can deliver augmented experiences with minimal mechanical overhead when transition costs are minimized: with sub-second onset times, even a three-second glimpse devotes the majority of the encounter to augmented content. This is not a compromise but a design choice---and one that appears to bring benefits in social comfort, shareability, and naturalness that sustained wear may not match.

Our findings also suggest that MR experience may not accumulate only linearly. The orbit behavior (Section~5.2) indicates that temporal wearing can support a \textit{cumulative} model of experience, where meaning builds across multiple brief encounters rather than within a single extended session. This pattern resonates with how people experience many cultural artifacts---returning to a painting, re-reading a poem, revisiting a scenic view---and suggests that MR content designed for temporal wearing may prove engaging in ways that differ from content designed for sustained immersion.

\subsection{What the Baseline Comparison Revealed}

The session-wear baseline in Phase~1 serves a critical function: it anchors our claims about temporal wearing in a direct comparison rather than relying solely on self-reported contrasts with prior HMD experience. [PLACEHOLDER: Discuss key findings from the baseline comparison. Expected topics:]

The comparison suggests several patterns worth noting. First, transition cost: the session-wear condition required [PLACEHOLDER: X] seconds for donning (strap adjustment, IPD calibration, orientation) compared to [PLACEHOLDER: <2] seconds for all temporal wearing conditions. This [PLACEHOLDER: order-of-magnitude] difference appears to be not just quantitative but qualitative---participants described the session-wear don as a ``setup'' process and the temporal wearing don as a ``gesture.''

Second, social comfort: participants reported [PLACEHOLDER: lower] social comfort in the session-wear condition, even in the controlled lab setting. Bystander companions reported [PLACEHOLDER: findings about watching someone in an HMD vs. watching someone peek through a handheld viewer]. These findings are consistent with the social acceptability literature~\cite{koelle2020social} but extend it by providing a direct within-subjects comparison between temporal and session wearing of the \textit{same content}.

Third, cognitive load: the NASA-TLX comparison revealed [PLACEHOLDER: findings]. We note that [PLACEHOLDER: any surprising findings, e.g., experience quality ratings were comparable despite the much briefer exposure in temporal wearing conditions, suggesting that temporal wearing's information density per second may compensate for shorter total exposure].

These baseline findings should be interpreted with caution. The lab setting is artificial, the session-wear condition used a specific HMD model, and the content was designed for temporal wearing (brief, self-contained) rather than for sustained session use. A fairer comparison might use content optimized for each wearing mode. Our findings indicate that temporal wearing differs from session wear on measurable dimensions, but the relative advantages of each mode likely depend on context, content, and user goals.

\subsection{Design Principles for Temporal Wearing}

Drawing on our design process and empirical findings from both study phases, we propose six design principles for temporal wearing devices. Each principle is grounded in specific findings; we note the supporting evidence and moderate our claims accordingly.

\subsubsection{P1: Design for Discovery.}
The device should invite curiosity before any interaction begins. Our deployment data suggests that the \textit{attract} phase of the glimpse cycle (Section~3.5) plays an important role: visitors who noticed the device---through its physical presence, through signage, or through observing another visitor's reaction---appeared more likely to engage than those who passed by unaware. The honeypot effect~\cite{wouters2016honeypot}---where one person's use attracts others---emerged as a discovery mechanism in [PLACEHOLDER: \%] of first interactions (Phase~2 video coding).

Design implications: Temporal wearing devices should be physically present and visually legible. Their form should communicate ``look through me'' through affordances~\cite{norman2013design}: lens-like apertures, viewfinder shapes, or transparent optics that hint at something beyond.

\subsubsection{P2: Instant-On.}
The transition from not-wearing to augmented perception should take less than two seconds. Our Phase~1 data indicates that transition cost is a primary mechanical correlate of re-engagement frequency: participants re-engaged more frequently with prototypes that offered near-instantaneous onset. Phase~2 deployment data is consistent with this pattern ($\rho$ = [PLACEHOLDER], $p$ = [PLACEHOLDER]).

Design implications: Displays should be always-on or wake from sleep in under 100ms. Tracking should be pre-initialized or use inside-out methods that do not require calibration. A useful heuristic: if the user can complete the don gesture faster than the device can deliver content, the device is too slow.

\subsubsection{P3: Graceful Doff.}
Exiting the augmented experience should feel natural, not abrupt or punishing. Interview participants in both phases contrasted temporal wearing with VR experiences where ``taking off the headset'' felt like ``being ejected from'' the experience. Temporal wearing doffing should feel like lowering binoculars---a natural return to unmediated perception.

Design implications: Avoid content that requires the user to complete an action before doffing. Design content that is complete at any moment of exit---every frame should be a valid stopping point. Optical see-through displays inherently support this by allowing the real world to persist behind the augmented layer.

\subsubsection{P4: Social Legibility.}
Bystanders should be able to perceive what the user is doing, why, and that it will be brief. Our social comfort data from both phases suggests that temporal wearing's social advantage depends on legibility: participants felt comfortable using the devices when the interaction ``looked normal'' and was visibly brief. The Phase~1 baseline comparison indicates that session-wear's lower social comfort scores may stem partly from reduced legibility---bystanders found it harder to interpret what the HMD user was doing and when they would finish.

Design implications: The donning gesture should reference familiar actions: raising binoculars, looking through a viewfinder, leaning toward a telescope. The device's form should communicate its purpose even to people who have never seen it before.

\subsubsection{P5: Invitation to Return.}
Content should reward re-engagement rather than penalize brevity. The orbit behavior (Phase~2) suggests that temporal wearing naturally invites return visits---but our interview theme of ``wanting more (in the right way)'' (Section~5.2) captures a critical distinction: temporal wearing should create anticipatory desire (``I want to come back'') rather than deprivation (``I wish it lasted longer'').

Design implications: Content should vary across encounters---through spatial variation, temporal variation, progressive revelation, or stochastic elements. The ideal is a ``generous preview'': each glimpse is complete and satisfying in itself, yet hints at more to be discovered.

\subsubsection{P6: Shared Glimpsing.}
Temporal wearing devices should enable, not obstruct, social interaction around the moment of augmented perception. The narrator, hand-off, and social permission behaviors (Phase~2) suggest that temporal wearing can produce shared MR experiences without requiring multiple synchronized devices. This shareability appears to be one of temporal wearing's most distinctive properties compared to session wear.

Design implications: Design for passability: handheld devices should be easy to pass between people. Design for narration: the user should be able to speak while wearing the device and should maintain awareness of companions. Design for spectatorship: bystanders should find it interesting to watch someone use the device.

\subsection{Temporal Wearing Beyond Mixed Reality}

While we have developed temporal wearing in the context of MR, the design dimension may extend to other domains of wearable and ubiquitous computing where brief, bounded device encounters offer advantages over sustained use.

\paragraph{Internet of Things.}
Many interactions with smart home devices are temporal: a brief check of a thermostat display, a momentary glance at an energy monitor. Designing these interactions explicitly with temporal wearing principles---instant-on, graceful exit, invitation to return---could improve their integration into domestic routines.

\paragraph{Health Monitoring.}
For many health-conscious individuals, a brief daily check---temporal wearing of a diagnostic device---may be more appropriate than continuous monitoring. Designing health devices for temporal wearing could reduce the burden and stigma of health monitoring while still providing actionable data.

\paragraph{Accessibility.}
Some users of assistive technology prefer intermittent use---engaging with assistive technology only in specific contexts~\cite{profita2016acceptability}. Temporal wearing offers a design framework for assistive devices that support this pattern: rapid onset when needed, graceful doff when not, and social legibility that reduces stigma.

\subsection{Limitations and Future Work}

Our work has several limitations that suggest directions for future research.

\paragraph{Lab Study Limitations.}
The Phase~1 lab study (N=20) provides controlled comparisons but in an artificial setting. The session-wear baseline used a single HMD model with content designed for temporal wearing, which may disadvantage the session-wear condition. Future work should compare temporal and session wearing with content optimized for each mode. The sample size, while adequate for within-subjects comparisons, limits generalizability; larger studies with demographic stratification would strengthen these findings.

\paragraph{Single Deployment Context.}
Our Phase~2 findings are drawn from a single deployment at [VENUE PLACEHOLDER]. While we selected this venue for its representative qualities, temporal wearing may manifest differently in other contexts---outdoor environments, retail settings, educational institutions, or domestic spaces. Future work should examine temporal wearing across diverse deployment contexts.

\paragraph{Content Effects.}
All prototypes presented the same augmented content, allowing us to isolate form factor effects. However, content type likely interacts with temporal wearing behavior: narrative content may encourage longer glimpses, while abstract visual content may support faster cycling. Systematically varying content alongside form factor would clarify these interactions.

\paragraph{Novelty Effect.}
Both study phases captured participants' first encounters with temporal wearing prototypes. The novelty of the devices may have inflated engagement metrics and positive responses. Longitudinal deployments---measuring how temporal wearing behaviors evolve over repeated visits across days or weeks---would address whether the patterns we observed are durable. We note, however, that the historical persistence of temporal wearing devices (coin-operated viewfinders, View-Masters) suggests that the practice itself may be robust to novelty decay, even if specific devices lose appeal.

\paragraph{Head-Mounted Temporal Wearing.}
Our prototype set did not include a head-mounted temporal wearing device. While head-mounted devices inherently incur higher transition costs, it is possible to design lightweight masks or flip-up visors that support temporal wearing from a head-mounted position. Future work should explore whether head-mounted form factors can achieve the sub-two-second transition costs that our findings suggest are important for temporal wearing.

\paragraph{Longitudinal Content Design.}
We proposed a layered revelation content model but did not empirically test it. Designing and deploying content specifically structured for cumulative glimpses, and measuring how visitor engagement evolves across multiple return visits, is a natural next step.

% Section 7: Conclusion
% Paper: "Just a Glimpse: Temporal Wearing as a Design Paradigm for Transient Mixed Reality"

\section{Conclusion}

This paper has introduced \textit{temporal wearing} as a distinct interaction paradigm for mixed reality---one in which the act of donning and doffing a device constitutes the interaction itself, with encounters lasting seconds rather than sessions. We formalized this paradigm through a three-part wearing duration taxonomy that distinguishes temporal wear from permanent and session wear, and articulated a design space defined by body coupling, transition cost, and social legibility. Through six prototypes spanning handheld, neck-hung, and stationary form factors, we instantiated diverse positions within this design space, and through a field deployment at [VENUE PLACEHOLDER], we demonstrated that temporal wearing produces interaction patterns---double-takes, hand-offs, narration, and orbiting---that are qualitatively distinct from session-based MR use.

Our contributions are fourfold. First, the \textbf{temporal wearing concept and duration taxonomy} (C1) name a centuries-old but previously untheorized practice, identifying wearing duration as a primary design variable that existing wearable computing frameworks have overlooked. Second, the \textbf{design space} (C2) provides a structured vocabulary for reasoning about temporal wearing devices and situating them relative to existing MR technology. Third, our \textbf{empirical findings} (C3) establish that transition cost shapes re-engagement frequency, that temporal wearing enables socially shared MR without multi-user systems, and that visitors develop distinctive behavioral patterns around brief, repeated encounters. Fourth, our \textbf{six design principles} (C4)---Design for Discovery, Instant-On, Graceful Doff, Social Legibility, Invitation to Return, and Shared Glimpsing---offer actionable guidance for designers working with transient MR encounters.

More broadly, temporal wearing challenges the implicit assumption that longer wearing is better wearing. The history of optical instruments---lorgnettes, stereoscopes, viewfinders, View-Masters---demonstrates that brief, purposeful encounters with augmented perception have been valued for centuries. The contemporary pursuit of all-day, always-on wearable computing is a valid design direction, but it is not the only one. The full spectrum of wearing durations---from seconds to days---deserves design attention, and the brief end of that spectrum has been neglected.

We envision a future in which wearable computing embraces this full spectrum. Alongside devices designed to be worn and forgotten, there should be devices designed to be picked up and put down---devices that invite glimpses, reward returns, and support the social sharing of augmented perception. A visitor raises a viewer to her eyes. For a few seconds, the world is more than it appears. She lowers it, turns to her companion, and says: ``You have to see this.'' That moment---brief, shared, complete---is what temporal wearing is for.


% All sections included

\bibliography{references}

% Appendix
% Paper: "Just a Glimpse: Temporal Wearing as a Design Paradigm for Transient Mixed Reality"

\appendix

\section{Prototype Descriptions}
\label{appendix:prototypes}

This appendix provides detailed descriptions of all six temporal wearing prototypes. Each description includes the design rationale, technical specifications, form factor characteristics, and transition profile.

\subsection{Prototype A: [PLACEHOLDER Name] (Handheld)}

\begin{figure}[h]
    \centering
    \includegraphics[width=0.8\columnwidth]{figures/prototype-a-placeholder}
    \caption{Prototype A: [PLACEHOLDER]. A handheld monocular/binocular viewer designed for single-hand operation. [PLACEHOLDER: brief description of appearance and use gesture.]}
    \label{fig:prototype-a}
\end{figure}

\paragraph{Design Rationale.}
[PLACEHOLDER: 2--3 sentences on why this form factor was chosen, what design space position it occupies, and what historical precedent it draws from (e.g., lorgnette, magnifying glass).]

\paragraph{Technical Specifications.}
\begin{itemize}
    \item Display: [PLACEHOLDER: type, resolution, FOV]
    \item Tracking: [PLACEHOLDER: method]
    \item Weight: [PLACEHOLDER] g
    \item Dimensions: [PLACEHOLDER] mm
    \item Power: [PLACEHOLDER: battery type, life]
    \item Onset time: [PLACEHOLDER] s (measured)
    \item Offset time: [PLACEHOLDER] s (measured)
\end{itemize}

\paragraph{Transition Profile.}
[PLACEHOLDER: Description of the don/doff gesture, what makes it fast, any design features that support instant-on.]

\subsection{Prototype B: [PLACEHOLDER Name] (Handheld)}

\begin{figure}[h]
    \centering
    \includegraphics[width=0.8\columnwidth]{figures/prototype-b-placeholder}
    \caption{Prototype B: [PLACEHOLDER]. [PLACEHOLDER: brief description.]}
    \label{fig:prototype-b}
\end{figure}

[PLACEHOLDER: Same structure as Prototype A --- design rationale, technical specifications, transition profile.]

\subsection{Prototype C: [PLACEHOLDER Name] (Neck-hung)}

\begin{figure}[h]
    \centering
    \includegraphics[width=0.8\columnwidth]{figures/prototype-c-placeholder}
    \caption{Prototype C: [PLACEHOLDER]. A neck-hung viewer suspended from a lanyard, resting on the chest when not in active use. [PLACEHOLDER: brief description.]}
    \label{fig:prototype-c}
\end{figure}

\paragraph{Design Rationale.}
[PLACEHOLDER: Rationale for neck-hung form factor. How passive body coupling enables rapid transition from rest to active use. Historical precedent: press photographer cameras, tourist binoculars worn on neck straps.]

\paragraph{Technical Specifications.}
\begin{itemize}
    \item Display: [PLACEHOLDER]
    \item Tracking: [PLACEHOLDER]
    \item Weight: [PLACEHOLDER] g
    \item Lanyard: [PLACEHOLDER: material, length, adjustment]
    \item Onset time: [PLACEHOLDER] s
    \item Offset time: [PLACEHOLDER] s
\end{itemize}

\paragraph{Transition Profile.}
[PLACEHOLDER: Description of the raise-from-chest gesture, how gravity assists offset.]

\subsection{Prototype D: [PLACEHOLDER Name] (Neck-hung)}

\begin{figure}[h]
    \centering
    \includegraphics[width=0.8\columnwidth]{figures/prototype-d-placeholder}
    \caption{Prototype D: [PLACEHOLDER]. [PLACEHOLDER: brief description.]}
    \label{fig:prototype-d}
\end{figure}

[PLACEHOLDER: Same structure --- design rationale, technical specifications, transition profile.]

\subsection{Prototype E: [PLACEHOLDER Name] (Stationary)}

\begin{figure}[h]
    \centering
    \includegraphics[width=0.8\columnwidth]{figures/prototype-e-placeholder}
    \caption{Prototype E: [PLACEHOLDER]. A plinth-mounted or wall-mounted viewfinder that users lean into. [PLACEHOLDER: brief description.]}
    \label{fig:prototype-e}
\end{figure}

\paragraph{Design Rationale.}
[PLACEHOLDER: Rationale for stationary form factor. The body approaches the device rather than vice versa. Historical precedent: coin-operated viewfinders, museum stereoscopes, peep shows. How fixed placement maximizes social legibility and eliminates device management.]

\paragraph{Technical Specifications.}
\begin{itemize}
    \item Display: [PLACEHOLDER]
    \item Tracking: [PLACEHOLDER]
    \item Mounting: [PLACEHOLDER: plinth/wall/pedestal, height, adjustability]
    \item Viewing aperture: [PLACEHOLDER: dimensions, eye relief]
    \item Onset time: [PLACEHOLDER] s
    \item Offset time: [PLACEHOLDER] s
\end{itemize}

\paragraph{Transition Profile.}
[PLACEHOLDER: Description of the lean-in gesture, proximity-triggered activation, natural withdrawal.]

\subsection{Prototype F: [PLACEHOLDER Name] (Stationary)}

\begin{figure}[h]
    \centering
    \includegraphics[width=0.8\columnwidth]{figures/prototype-f-placeholder}
    \caption{Prototype F: [PLACEHOLDER]. [PLACEHOLDER: brief description.]}
    \label{fig:prototype-f}
\end{figure}

[PLACEHOLDER: Same structure --- design rationale, technical specifications, transition profile.]


\section{Interview Protocol}
\label{appendix:interview}

Semi-structured exit interview guide used with visitors following interaction with temporal wearing prototypes. Interviews lasted 10--15 minutes and were audio-recorded with consent.

\subsection{Opening}
\begin{enumerate}
    \item Thank you for agreeing to speak with us. Can you tell me a bit about your visit today? Who are you here with?
    \item I noticed you used [prototype name/description]. Can you walk me through what happened, from when you first noticed it?
\end{enumerate}

\subsection{Discovery and Approach}
\begin{enumerate}[resume]
    \item What first caught your attention about the device?
    \item What did you think it was before you used it?
    \item Was there anything that made you hesitate before picking it up / leaning in?
\end{enumerate}

\subsection{The Experience}
\begin{enumerate}[resume]
    \item Can you describe what you saw when you looked through it?
    \item How long do you think you looked through it the first time?
    \item Did you look through it more than once? [If yes:] What made you want to look again?
    \item How did it feel to put it down / step away?
\end{enumerate}

\subsection{Social Dimensions}
\begin{enumerate}[resume]
    \item Were you aware of other people around you while you were using it?
    \item Did you feel self-conscious at all? [Probe: why/why not?]
    \item Did you share the experience with anyone? [If yes:] How?
    \item [If in group:] Did your companion(s) also try it? How did that come about?
\end{enumerate}

\subsection{Comparison and Reflection}
\begin{enumerate}[resume]
    \item Have you used VR or AR devices before? [If yes:] How did this compare?
    \item If this device were available at other places you visit---museums, parks, stores---would you use it?
    \item Is there anything else you'd like to share about the experience?
\end{enumerate}


\section{Survey Instruments}
\label{appendix:survey}

\subsection{User Social Comfort Survey}

Administered immediately after interaction. All items rated on a 7-point Likert scale (1 = Strongly Disagree, 7 = Strongly Agree).

\begin{enumerate}
    \item I felt comfortable using this device in this setting.
    \item I was aware of other people watching me while I used the device.
    \item I felt self-conscious while using the device. [R]
    \item The act of picking up / leaning into the device felt natural.
    \item I would feel comfortable using this device in a different public setting (e.g., a park, a shop).
    \item I would use this device again if I encountered it.
    \item The device was easy to start using (no instructions needed).
    \item I felt socially disconnected from my companions while using the device. [R]
\end{enumerate}

Items marked [R] are reverse-coded.

\subsection{Bystander Comfort Survey}

Administered to companions or nearby visitors who observed but did not use the device.

\begin{enumerate}
    \item I could tell what the person was doing with the device.
    \item The person's use of the device seemed natural in this setting.
    \item I felt uncomfortable watching the person use the device. [R]
    \item I understood that the interaction would be brief.
    \item I felt the person was still ``present'' and available during their use of the device.
    \item Seeing the person use the device made me want to try it myself.
\end{enumerate}


\section{Video Coding Scheme}
\label{appendix:coding}

Adapted from Brignull and Rogers~\cite{brignull2003enticing} and extended for temporal wearing contexts. Each interaction episode (defined as a continuous sequence from first approach to final departure from the device vicinity) was coded on the following dimensions.

\subsection{Discovery Phase}
\begin{itemize}
    \item \textbf{Discovery trigger}: Self-initiated (noticed device independently) / Social (prompted by companion or observed another user) / Signage (read label or sign) / Staff (directed by venue staff)
    \item \textbf{Approach trajectory}: Direct (walked straight to device) / Circuitous (approached indirectly, pausing at other exhibits) / Return (came back after initial pass)
    \item \textbf{Hesitation}: None / Brief pause ($<$3s) / Extended pause ($>$3s) / Abandoned approach
\end{itemize}

\subsection{Donning Phase}
\begin{itemize}
    \item \textbf{Donning gesture}: One-hand lift / Two-hand lift / Lean-in / Raise from chest / Other
    \item \textbf{Onset confidence}: Immediate (no adjustment) / Minor adjustment (1 repositioning) / Multiple adjustments / Required assistance
    \item \textbf{Onset time}: Measured in seconds from first contact to stable viewing position
\end{itemize}

\subsection{Glimpse Phase}
\begin{itemize}
    \item \textbf{Glimpse duration}: Measured in seconds
    \item \textbf{Viewing behavior}: Stationary gaze / Scanning (moving device or head) / Searching (active exploration) / Showing companion
    \item \textbf{Verbal behavior during glimpse}: Silent / Exclamation / Narration to companion / Question to companion
    \item \textbf{Body posture}: Upright / Leaning / Crouching / Other
\end{itemize}

\subsection{Doffing Phase}
\begin{itemize}
    \item \textbf{Doffing gesture}: Lower to surface / Lower to chest / Hand to companion / Step back / Other
    \item \textbf{Offset time}: Measured in seconds
    \item \textbf{Facial expression at doff}: Neutral / Smile / Surprise / Puzzlement / Other
\end{itemize}

\subsection{Post-Doff Phase}
\begin{itemize}
    \item \textbf{Immediate action}: Verbal sharing / Gestural sharing (pointing) / Re-engagement (picks up again) / Moves on / Waits for companion
    \item \textbf{Social interaction}: None / Narration / Discussion / Hand-off / Joint re-engagement
    \item \textbf{Return}: No return / Return within 1 minute / Return within 5 minutes / Return after 5+ minutes
\end{itemize}

\subsection{Temporal Wearing Behaviors}
\begin{itemize}
    \item \textbf{Double-take}: Quick initial glimpse ($<$3s) followed by re-engagement within 10s
    \item \textbf{Hand-off}: Device passed directly to companion within 15s of doffing
    \item \textbf{Narrator}: User verbally describes augmented content to companion during glimpse
    \item \textbf{Orbit}: User returns to device after visiting other parts of venue ($>$1 minute interval)
\end{itemize}


\end{document}
